\epigraph{Civilization advances by extending the number of important operations which we can perform without thinking about them}{Alfred North Whitehead, An Introduction to Mathematics}

\lettrine{A}n optimizing compiler must preserve the semantics of the source program in the generated target program,
and moreover, must do so while optimizing the target program to run as efficiently as possible.
This makes implementing production-grade compilers complex, and proving them correct even more so.
For example consider the LLVM compiler framework. Its peephole optimizer,
which scans for a small window of instructions and replaces them with a more efficient equivalent sequence of instructions, is conceptually a simple find-and-replace.
However, in terms of the actual implementation, the peephole rewrite is enormous, and is about 10\% of the entire LLVM codebase.
Moreover, proving this correct requires one to prove the correctness of each individual rewrite, of which there are thousands,
and each rewrite is non-trivial in its own right.
In this thesis, I focus on building techniques to allow us to prove such peephole rewrites correct in a proof assistant.

Toward verified peephole rewriting, I first formalize a core calculus for SSA-based compiler intermediate representations,
and prove a generic peephole rewriter correct.
This demonstrates a framework for verifying SSA-based peephole rewriting in a proof assistant.
I also prove a lifting theorem,
which says that proving a rewrite correct in isolation is sufficient to prove it correct in the context of a whole program.
Overall, this discharges the structural half of the problem of verifying a peephole rewriter.
However, the other half, of actually proving the rewrite correct, still remains.
Effective proof automation is necessary to make this approach scalable to the thousands of rewrites in LLVM.

In the second half of the thesis, I introduce new, formally verified decision procedures for proving such rewrites correct automatically.
Our new decision procedures can prove the correctness of parametric bitvectors and floating point.
Moreover, by mechanizing the decision procedures in a proof assistant, we are assured of their correctness.
Algorithmically, our decision procedures are based on bitblasting,
which translates problems from a given theory into propositional logic,
and then solves the propositional formula using a SAT solver.
For the verification pipeline, I build on \bvdecide{}, Lean's formally verified fixed-width bitblasting infrastructure.

%Perhaps counterintuitively, the proof obligations a compiler hands us are not immediately bitblastable.
%For example, a rewrite has to be correct at \emph{every} register width, whereas a bitblaster wants one fixed width,
%and the standard remedy (used, for example, by the Alive peephole rewrite verifier) is to enumerate the widths and bitblast each one separately.
First, I create new state-of-the-art algorithms for parametric bitvector theory,
which is a generalization of fixed-width bitvectors, where the width is a symbolic variable rather than a concrete number.
This is necessary for verifying compiler rewrites, as a rewrite must be correct at all bitwidths.
Until now, the standard technique, as good decision procedures for parametric bitvectors are unavailable,
has been to enumerate the widths and bitblast each one separately.
I show that this is unnecessary.
Predicates that must hold at all widths reduce to reachability in a finite automaton, which is decided by standard algorithms from model checking.
Next, we also show that any formula over $n$ symbolic widths reduces to an equisatisfiable single-width formula.
This improves from an exponential number of solver queries for a \emph{single} query,
and provides a new sound and complete fragment of multi-width theory.

Second, I build the first formally verified floating-point bitblaster.
Floating-point semantics is intricate, and machine learning is pushing exotic formats such as E4M3 and E5M2 into production hardware,
precisely where solvers are least tested.
Toward trustworthy floating-point verification, I mechanize the SMT-LIB \qffp{} theory,
and build verified bitblasting circuits for almost all of its operations.

In summary, our width-generic solvers solve more all-width problems than the unverified state of the art,
and our floating-point bitblaster, while slower than Bitwuzla,
offers in exchange end-to-end guarantees down to Lean's kernel.
I demonstrate that foundational proof and push-button automation need not be at odds,
and that the reductions developed here lay the groundwork for effective verification of peephole rewriting at scale.
